\chapter{System Description}

\section{Overview and Design Rationale}

The experimental system is a web-based apparatus for studying touch-driven list scrolling under a configurable transfer function. Its purpose is to present participants with a controlled search-and-select task on a long, continuously scrollable list, and to adapt the scrolling dynamics to each participant over the course of a session by means of an online optimization loop.

The guiding architectural principle is a strict separation of concerns across three layers. A React single-page application presentation layer renders the list and takes full programmatic control of touch input and scrolling. A managed GraphQL API backed by a schemaless document store persistence layer records participants and their completed trial blocks. A serverless machine-learning backend optimization layer consumes the recorded data on demand and proposes the next transfer-function parameter set.

A second, central design decision concerns the scrolling behaviour itself. Rather than relying on the browser's native momentum scrolling or on a framework's default scroll view, the post-touch (fling) motion is governed by an explicit, parameterized transfer function that was derived directly from the source code of Android's OverScroller class. Exposing the constants that are otherwise fixed in that source as tunable parameters turns regular preconfigured scrolling into a well-defined numerical search space, which the optimization layer can then adapt per participant. This coupling between a physically motivated interaction model and a data-driven optimizer is the conceptual core of the system.

\section{System Architecture}

The system is event-driven from end to end. Raw interaction data are recorded by the frontend, persisted in blocks of ten trials, evaluated on demand by a Gaussian-Process-based optimizer and the resulting new transfer-function parameters are made available to the participant for the subsequent block. The following subsections describe the four building blocks of this architecture.

\subsection{Frontend}

The frontend is a React single-page application written in plain JavaScript. Its static build artifacts (HTML, JavaScript, CSS) are hosted in an Amazon S3 bucket and delivered to participants through an Amazon CloudFront content-delivery distribution, which also terminates HTTPS. The frontend is responsible for rendering the list, running the custom touch-input and scrolling pipeline, evaluating the transfer function locally in real time, and collecting per-trial measurements that are submitted to the backend one block at a time.

\subsection{Backend and Data Flow}

The application backend is built around AWS AppSync, a managed GraphQL API service, which exposes queries and mutations for participant management, for storing completed trial blocks, and for requesting a new transfer-function parameter set. Amazon DynamoDB serves as the persistent, schemaless data store behind this API: each participant is represented as a single item holding demographic fields (name, e-mail, birth date, smartphone ownership, daily screen time), the currently active transfer-function parameter set, the next parameter set once available, and the complete history of completed trial blocks stored as a nested JSON attribute. Because a participant's entire state lives in one item, reading and appending trial data require no joins and remain a single-item operation.

\subsection{Optimization Backend}

The optimization backend is invoked synchronously as part of the GraphQL request flow. When a participant completes a block of trials, the frontend calls the mutation that requests the next parameter set. The AppSync resolver forwards this request to an AWS Lambda function, which reads the participant's stored trial history from DynamoDB and combines it with the anonymized histories of the other participants recorded so far. The resulting payload is sent to an Amazon SageMaker inference endpoint that performs the parameter selection for the next block.

The inference service is implemented as a stateless optimization routine. On each request, the endpoint reconstructs a participant-aware training table from the incoming block-level observations, validates the data, and fits a fresh multi-objective Bayesian optimization model over the available parameter space. The model is conditioned on the participant identity as an additional input dimension, so that information from previous participants can inform the prediction for the currently active participant while still preserving participant-specific differences. Each of the ten target positions in a block contributes one objective, and the acquisition function selects the parameter set that best trades off average completion time and target-dependent variability.

The resulting candidate parameter set is returned to the Lambda function, which writes it back to DynamoDB as the participant's next parameter set. The frontend then retrieves this value once the participant explicitly confirms readiness for the next block. This design keeps the optimization loop closed over the experimental data while maintaining a clear separation between data collection, inference, and the subsequent trial block.

\subsection{Infrastructure as Code and Deployment}

All infrastructure, the S3 bucket and CloudFront distribution, the AppSync API and its resolvers, the DynamoDB table, the Lambda function, and the SageMaker endpoint, is defined declaratively as code with Terraform and deployed through a GitHub Actions CI/CD pipeline. On every push to the main branch, the pipeline authenticates against AWS via OpenID Connect, runs terraform apply to reconcile the deployed AWS infrastructure with its version-controlled definition, and then builds and publishes the React frontend. This keeps the presentation, persistence, and optimization layers independently deployable while guaranteeing that the configuration used during data collection is fully reproducible from the repository.

\section{Touch Input and Scrolling Pipeline}

To obtain full and reproducible control over the scrolling dynamics, the frontend does not use the browser's native scrolling. Instead, it processes raw touch events and drives the list position programmatically, so that the fling behaviour is determined solely by the transfer function and not by browser or device-specific momentum implementations.

\subsection{List and Scroll Container}

The user interface presents a single continuously scrollable list of $N = 330$ entries, numbered consecutively. The list is not paginated in the sense of discrete page breaks or page-navigation controls. Instead, the height of each list entry is computed dynamically from the available viewport height so that exactly 15 entries are visible within the scroll container at any given time. This keeps the spatial density of the list, and therefore the visual arrangement participants perceive, consistent across devices and across trials.

\subsection{Gesture Processing}

The browser's and React's default scrolling behaviour is disabled for the list container and replaced by a custom pipeline that processes the native onTouchStart, onTouchMove, onTouchEnd, and onTouchCancel events directly. During a touch gesture, every touch-move sample is time-stamped and stored in a rolling buffer bounded by a $100$\,ms window and at most 20 samples. For real-time drag feedback the list follows the finger through an exponentially smoothed instantaneous velocity, while for launching a fling the release velocity is estimated from the buffered samples by a quadratic least-squares fit, taking the slope of the fitted curve at the moment of release as the launch velocity. This is more robust than a two-point derivative and, unlike a straight-line fit, follows the acceleration typically present at the end of a fling. If the estimated velocity exceeds a fixed fling threshold, it is clamped and passed into the fling model. Otherwise the list stops at its current position, after being clamped back into the valid scroll bounds if the participant had dragged past the list's edges.

\subsection{Overscroll and Bounds Handling}

Android's native scrolling system treats the list boundary as a special state. When a drag or fling would push the content beyond the valid range, the system allows a limited overscroll excursion and applies a resistance effect so that the visible motion decays smoothly near the boundary instead of abruptly stopping. Once the user releases the finger, or once the fling reaches the end of the valid range, the content is pulled back to the nearest legal position. This behavior is part of the general edge-response logic of the platform and is distinct from the actual fling trajectory itself: the ballistic motion is governed by the deceleration model, whereas the boundary handling is a separate mechanism that prevents the list from leaving its legal bounds.

In this work, the same principle is implemented in a deliberately simplified form. The scroll position is always kept inside the valid list range, and any overshoot beyond the top or bottom is clamped and damped before being visualized. Unlike the full Android implementation, the project does not model the entire native edge state machine or its spring-back dynamics in detail; instead, it separates the concerns cleanly: the transfer function governs the in-bounds fling motion, while the bounds logic acts as a fixed, parameter-independent guard that stops the list at the valid edge and restores it to a legal position when necessary.

\section{Transfer Function}

The transfer function governing fling motions is a physically motivated model derived directly from Android's OverScroller spline-deceleration algorithm.

\subsection{Fling Phase}

The fling phase is the normal, in-bounds momentum motion that follows a flick gesture. It is built on a strict separation between the scale of a fling (how far and how long it travels, computed from the launch velocity and a few constants) and its shape (the fixed, dimensionless temporal profile with which that distance is covered. Playback then simply stretches the shape onto the computed scale).

\subsubsection{Scale: Physical Coefficient, Fling Distance, and Duration}

The scale converts the launch velocity $v_0$ into a total travel distance $D$ and a total duration $T$. It relies on a device-specific physical coefficient
\begin{equation}
  C \;=\; g \cdot 39.37 \cdot \mathrm{ppi} \cdot 0.84,
  \qquad
  \mathrm{ppi} \;=\; 160 \cdot \rho,
  \label{eq:sd-physical-coeff}
\end{equation}
and on a fixed dimensionless deceleration rate
\begin{equation}
  r \;=\; \frac{\ln 0.78}{\ln 0.9} \;\approx\; 2.3582 .
  \label{eq:sd-decel-rate}
\end{equation}
Given a fling friction coefficient $\mu$, the distance and the duration (the latter in milliseconds) follow in closed form as
\begin{align}
  D(v_0) &= \mu\,C \left( \frac{\beta\,\lvert v_0 \rvert}{\mu\,C} \right)^{\!\frac{r}{r-1}}, \label{eq:sd-fling-distance}\\[2pt]
  T(v_0) &= 1000 \left( \frac{\beta\,\lvert v_0 \rvert}{\mu\,C} \right)^{\!\frac{1}{r-1}}. \label{eq:sd-fling-duration}
\end{align}
The parameters are:
\begin{itemize}
  \item $v_0$, the launch velocity (px/s).
  \item $g = 9.80665\,\mathrm{m/s^2}$, Earth's gravitational acceleration.
  \item $\rho$, the display density, giving $\mathrm{ppi} = 160\rho$ pixels per inch, with $39.37$ the inch-per-metre conversion factor and $0.84$.
  \item $C$, the resulting device-specific physical coefficient of Equation~\eqref{eq:sd-physical-coeff}.
  \item $\mu$, the dimensionless fling friction, which scales both distance and duration; in the original \texttt{OverScroller} it is the fixed constant $\mu = 0.015$ (\texttt{ViewConfiguration.SCROLL\_FRICTION}).
  \item $r \approx 2.3582$, the fixed deceleration rate of Equation~\eqref{eq:sd-decel-rate}.
  \item $\beta = 0.35$, the inflexion factor.
  \item the factor $1000$ in $T(v_0)$, converting the duration from seconds to milliseconds.
\end{itemize}

Because $\tfrac{r}{r-1} > 1$ while $\tfrac{1}{r-1} < 1$, distance grows super-linearly and duration sub-linearly with $v_0$: a faster fling covers disproportionately more ground while spreading it over a comparatively shorter time.

\subsubsection{Spline Curve and Normalized Sampling}

The scale provides the signed distance $D_s$ =  and the duration $T$. So that a single, fixed shape curve can describe every fling regardless of how far or how long it travels, these two quantities are normalized. Elapsed time is measured relative to $T$ and travelled distance relative to $D_s$, placing the motion in the dimensionless unit square $[0,1]\times[0,1]$. On this normalized domain the motion is a fixed position-versus-time spline $p(\tau)$ relating normalized time $\tau \in [0,1]$ to normalized distance $p \in [0,1]$. It is defined parametrically by two cubic B\'ezier curves sharing a common parameter $x \in [0,1]$; with the start and end tension constants $S = 0.5$ and $E = 1.0$ and the control ordinates
\begin{equation}
  P_1 = S\,\beta = 0.175,
  \qquad
  P_2 = 1 - E\,(1-\beta) = 0.35,
  \label{eq:sd-control-points}
\end{equation}
the normalized time and distance are
\begin{align}
  \tau(x) &= 3x(1-x)\big[(1-x)P_1 + x P_2\big] + x^3, \label{eq:sd-spline-time}\\
  p(x)    &= 3x(1-x)\big[(1-x)S + x\big] + x^3. \label{eq:sd-spline-pos}
\end{align}
This curve is sampled once, ahead of any fling and independently of $D_s$ and $T$, into a lookup table of $N+1 = 101$ entries ($N = 100$) holding $p$ at the equally spaced normalized times $\tau = i/N$.

At run time this normalization is inverted to recover pixels. At elapsed time $t$ (ms) the duration $T$ normalizes time to $\tau = t/T \in [0,1]$, which indexes the table; the two surrounding entries yield the normalized position $p(\tau)$ and its local slope $p'(\tau)$ by linear interpolation. The distance $D_s$ then scales both quantities back into pixels:
\begin{align}
  x(t) &= x_{\text{start}} + p(\tau)\,D_s, \label{eq:sd-spline-position}\\[2pt]
  v(t) &= p'(\tau)\,\frac{D_s}{T}\cdot 1000. \label{eq:sd-spline-velocity}
\end{align}

The parameters are:
\begin{itemize}
  \item $D_s$ and $T$, the signed distance and duration: $T$ normalizes elapsed time via $\tau = t/T$, while $D_s$ scales the looked-up values back into pixels.
  \item $S = 0.5$ and $E = 1.0$, the start and end tension constants that define the shape curve.
  \item $\beta = 0.35$, the inflexion factor, which enters the control ordinates $P_1$ and $P_2$.
  \item $x$, the shared B\'ezier parameter used to trace the shape curve.
  \item $N = 100$, the number of tabulated intervals ($N+1 = 101$ samples).
  \item $\tau = t/T$, $p(\tau)$ and $p'(\tau)$, the normalized time, the tabulated position and its interpolated slope.
  \item $x_{\text{start}}$, the list position at fling launch. The factor $1000$ converts the per-millisecond velocity back to pixels per second.
\end{itemize}

\subsection{Parameterization and Search Space}

\section{Adaptive Parameter Optimization}

Trials are organized into blocks of 10, and a single transfer-function parameter set remains fixed for the duration of one block. After each block, the system may propose a new parameter set for the next block, gradually adapting the scrolling behaviour to the participant.

\subsection{Warm-up and Bootstrap Phase}

For the first three blocks of each participant, the machine-learning backend is not yet used to propose parameters. Instead, each parameter is independently sampled uniformly at random from its predefined bounds. This warm-up/bootstrap phase serves two purposes: it generates an initial, sufficiently diverse set of observations for the optimization model to learn from, and it avoids relying on an uninitialized or poorly informed model at the very beginning of a session, where no data are yet available.

\subsection{Participant-Aware Multi-Objective Bayesian Optimization}

Only after these first three blocks have been completed does the system switch to model-based parameter generation. At that point, the optimization backend treats the problem as pooled and multi-objective rather than as a single-participant regression problem. Whenever a new parameter set is requested, the Lambda function retrieves the most recent block histories for the current participant together with the anonymized block histories of the previously recorded participants, and forwards the combined payload to the SageMaker endpoint.

The request payload is prepared in a structured form that reflects the actual optimization pipeline. Each recorded trial contains the completion time \texttt{timeMs}, the target index \texttt{targetNumber}, the block index \texttt{blockIndex}, and the transfer-function parameters active for that block in \texttt{paperParams}. On the inference side, the records are normalized into a canonical block-level schema, invalid parameter values are replaced with the default transfer-function constants, and the admissible parameter bounds are enforced. The records are then grouped by block, the parameter vector is taken from the first valid entry in the block, and the ten target positions are collapsed to their mean completion time per target. This yields a per-block observation of the form
\begin{equation}
    [\text{scrollFriction},\ \text{decelerationRate},\ \text{inflexion}] \rightarrow [t_{30}, t_{60}, \ldots, t_{300}],
\end{equation}
with the participant-specific context represented by an ordinal task feature that is assigned from the pooled participant ordering in the request payload rather than by the raw participant identifier itself. The pooled training set is thus built from block observations across participants, while preserving a participant-specific task dimension in the input space.

Duplicate or near-identical rows are collapsed by rounding the parameter values to six decimal places and averaging their objective vectors, which prevents repeated evaluations from making the Gaussian-process covariance numerically unstable. The optimizer then fits a multi-objective Bayesian model using BoTorch's \texttt{qLogNoisyExpectedHypervolumeImprovement} acquisition function. A separate Gaussian process is trained for each of the ten target-position objectives, while the task feature is included as part of the joint input space and the negative completion times are used as the objectives to be maximized. This keeps the model participant-aware without requiring a separate model per participant: the optimizer can borrow statistical strength across users while still distinguishing the participant slot through the task feature.

The optimization model is persisted across requests in the active SageMaker worker. When a new block is submitted, the endpoint updates the stored training table from the current pooled block histories, refits the model, and searches the bounded parameter space with a quantized acquisition-based candidate grid. The search samples a fixed set of candidate points from the admissible parameter ranges, normalizes them to the unit cube, evaluates their acquisition values under the current model, and selects the parameter vector with the highest value. A smaller probe set is used to approximate the rank of the selected candidate relative to random alternatives, and the resulting response includes both the selected parameter set and lightweight diagnostics such as the acquisition value, the estimated rank, and the number of training observations used for fitting. If the pooled observation count is below the minimum threshold required for Bayesian optimization, the endpoint refuses to update the model and falls back to the surrounding Lambda logic rather than proposing an unstable candidate.

\subsection{Candidate Scoring and Selection}

At each invocation, the inference endpoint samples a candidate grid of 256 quantized parameter points from the bounded search space, normalizes the inputs to a unit cube, and evaluates them under the current multi-objective acquisition function. The search is not performed with a hand-written scalar objective, but with a hypervolume-based acquisition criterion that jointly trades off all ten target-position completion times. A small probe set is used to approximate the rank of the selected candidate relative to random alternatives, and the candidate with the highest acquisition value is chosen as the next parameter set. The endpoint also includes lightweight metadata in the response, such as the acquisition value, candidate rank estimate, and the number of pooled observations used for fitting.

If fewer than two completed blocks of pooled data are available for the requesting participant, or if the model fails to fit or evaluate, the backend falls back to returning the participant's current parameter set unchanged. This prevents the optimization loop from proposing unstable values before the training data set is informative enough to support a meaningful model update.

\subsection{SageMaker Inference Implementation and Deployment}

The optimization model is not persisted as a learned artifact. Instead, a fresh Bayesian optimization run is carried out on every invocation from the current request payload, which keeps the deployment stateless and guarantees that optimization always reflects the most recent experimental data. Concretely, the inference logic is implemented as a single Python module exposing the required SageMaker entry points \texttt{model\_fn}, \texttt{input\_fn}, \texttt{predict\_fn}, and \texttt{output\_fn}. In the current implementation, \texttt{model\_fn} returns an empty model dictionary, while the data preparation, multi-objective GP fitting, candidate search, and response packaging are all performed inside \texttt{predict\_fn}. This module is packaged together with the lightweight required runtime files into a gzip-compressed \texttt{model.tar.gz} using SageMaker's \texttt{code/} layout, uploaded to S3 by Terraform, and loaded into a prebuilt PyTorch inference container via the \texttt{SAGEMAKER\_PROGRAM} and \texttt{SAGEMAKER\_SUBMIT\_DIRECTORY} environment variables.

This is a materially different deployment from the earlier Scikit-learn GP implementation: the endpoint uses PyTorch and BoTorch, since the acquisition function and multi-objective optimization logic depend on the PyTorch backend. The request payload is exchanged as \texttt{application/json} and carries the current parameter set together with the pooled per-participant block histories; the response contains the newly chosen parameter vector and a concise diagnostic bundle describing the acquisition strategy and the amount of training data that informed the decision.

\section{Experimental Procedure}
\label{sec:procedure}

\subsection{Task Design}
\label{sec:task}

Participants were asked to repeatedly locate and select a specified target entry within the 330-entry list. Ten fixed target positions were used throughout the experiment: entries $30, 60, 90, \dots, 300$, i.e.\ every 30th position starting 30 entries from the top and ending 30 entries from the bottom of the list. This excludes the first and last 30 entries from ever being used as a target, so that every target position is reachable by scrolling in either direction and no target is trivially visible without scrolling.

For each block of trials -- corresponding to one fixed transfer-function parameter set -- all 10 target positions were presented exactly once, in a randomized order that was freshly shuffled at the start of every block. Before the first randomized block, participants completed a single untimed practice trial with a fixed target (list entry 180) in order to familiarize themselves with the touch-scrolling interaction and the location-feedback display; this practice trial is not included in the recorded/optimized trial data.

\subsection{Interface Elements and Progress Feedback}
\label{sec:interface}

A horizontal progress indicator above the list shows, as a row of 10 pill-shaped segments, how many of the 10 trials belonging to the currently active transfer function have already been completed. Directly above the list, a banner displays the number the participant is currently searching for. Two identical vertical location-feedback tracks are rendered on the left and right edge of the screen, each showing two markers: one fixed at the position of the current target (proportional to its index within the 330-entry list) and one continuously updated to reflect the participant's current scroll position (computed from the on-screen offset of the topmost rendered items). This lets participants continuously orient themselves relative to the target without having to keep the target itself in view while scrolling. While a trial is in progress, the target entry is visually distinguished from all other (otherwise identically styled) entries by a red border and light red background; once the participant has tapped the correct entry, it briefly switches to a green background before the next trial begins.

\subsection{Data Collection and Recorded Metrics}
\label{sec:data-collection}

In addition to the timing data used by the optimizer, several gesture-level quantities are computed in real time during each trial for on-screen feedback and threshold-logic purposes, including swipe direction, per-gesture distance and duration, average and maximum swipe speed, the number of direction reversals (``switchbacks''/clutches) within a trial, and whether and by how far the target was overshot. Of these, the total elapsed time until the correct entry is tapped (\texttt{timeMs}), the net scrolled distance during the trial (\texttt{scrollDistance}), and a completion timestamp are recorded and persisted to the backend per trial; these are the quantities used as the optimization objective described above.

Trial data belonging to one block are collected locally in the browser and only submitted to the backend once all 10 trials of a block are complete, at which point they are appended atomically to the participant's stored trial history together with the transfer-function parameter set that was active during that block. After a block has been submitted, the participant is shown a waiting screen until a new parameter set has been generated (either through random sampling during the warm-up phase, or through the Gaussian-Process optimizer thereafter); the next block of trials is then only started once the participant explicitly confirms via a button click, creating a clear separation between successive optimization blocks and ensuring that a new transfer function is only introduced once the backend has produced an updated, validated parameter set.
